\section{Methodology}
\label{sec:methodology}

This section formalises the integration of a Gromov-Wasserstein
regularizer into a conditional flow-matching objective. We first
recall the flow-matching loss
\citep{cheng_2025_mmaudio}, then define the four families of
intra-modal pairwise-distance matrices that feed the regularizer,
the entropic and fused-entropic Gromov--Wasserstein objectives, the
Sinkhorn solver used to evaluate them, and the schedule that
controls their contribution to training. 

As the backbone of our experiments, we rely entirely on MMAudio \cite{cheng_2025_mmaudio}. The choice of MMAudio as the the backbone are motivated by (i) its state-of-the-art generative performances (ii) its lack of geometric relational awareness and (iii) its domain-relevancy with the previous experiments discussed in our thesis. 
Throughout, we treat
MMAudio purely as a conditioning architecture: a means of
materialising a video-conditioned generative process whose
representations the GW loss can shape. Our claims only concern the contributions of regularization methods on the backbone, but we are argue that the same methodology can be applied to the loss of any conditional generative model. 

\subsection{Backbone objective: conditional flow matching}
\label{sec:methodology:fm}

\paragraph{From CNFs to flow matching.}
A continuous normalising flow (CNF) parameterises a generative model
as the time-1 map of an ordinary differential equation
$\dot x_t = v_\theta(t, x_t)$, transporting samples
$x_0 \sim p_0$ from a tractable source (typically
$p_0 = \mathcal{N}(0, I)$) to samples $x_1 \sim p_1$ from the data
distribution. Classical maximum-likelihood training of CNFs
requires backpropagating through an ODE solver, which is
prohibitively expensive at the scales of interest here. Flow
matching \citep{lipman_2023_flow}, developed concurrently as
*stochastic interpolants*
\citep{albergo_2023_building, albergo_2023_stochastic}, sidesteps
this cost: rather than maximising the likelihood induced by the
ODE, one fixes a *conditional probability path*
$p_t(x \mid x_1)$ between $p_0$ and $p_1$, derives its target
velocity field $u_t(x \mid x_1)$ in closed form, and trains
$v_\theta$ by direct $L_2$ regression onto the target. Training is
then *simulation-free*: each gradient step only requires sampling
$(t, x_0, x_1)$ and a single forward pass through $v_\theta$, with
no SDE or ODE integration in the loop. Inference still solves an
ODE, but only at sampling time.

\paragraph{The OT / linear path.}
Among admissible probability paths, the one induced by
optimal-transport displacement interpolation between $p_0$ and a
single data point $x_1$ has a particularly simple form: the
deterministic linear interpolation
\begin{equation}
  x_t \;=\; t\,x_1 + (1-t)\,x_0,
  \qquad x_0 \sim \mathcal{N}(0, I), \;\; t \sim p_t,
  \label{eq:fm_path}
\end{equation}
with the constant target velocity
$u(x_t \mid x_0, x_1) = x_1 - x_0$. This linear path is the limit
of the Gaussian-path family at zero variance and underlies the
\emph{rectified flow} formulation of \citet{liu_2023_rectified};
its straightness makes the learned ODE accurately integrable in
few Euler steps. \citet{tong_2024_improving} extend the recipe to
non-Gaussian sources by coupling $(x_0, x_1)$ with a minibatch OT
plan, recovering the same loss with an empirically straighter
flow. We adopt the simplest variant — Gaussian source, independent
coupling, linear path — matching the rectified-flow choice of the
backbone we build on.

\paragraph{Conditional flow matching for V2A.}
For video-conditioned audio generation, let
$x_1 \in \R^{N \times d_a}$ denote a target audio latent of $N$
tokens of dimension $d_a$ (in our default configuration $N=250$ at
$31.25$ fps and $d_a=20$, the VAE latent space inherited from
MMAudio's S-16kHz model). Let
$\mathbf{C} = (F_v, F_{\mathrm{syn}}, F_t)$ collect the conditioning
features extracted by frozen encoders: CLIP visual features $F_v$,
Synchformer features $F_{\mathrm{syn}}$, and CLIP text features
$F_t$. Conditional flow matching learns a time-dependent velocity
field
$v_\theta(t, \mathbf{C}, x_t) : [0,1] \times \R^{|\mathbf{C}|}
\times \R^{N\times d_a} \to \R^{N \times d_a}$
that pushes the Gaussian prior toward the data distribution along
the linear path \eqref{eq:fm_path}, conditioned on $\mathbf{C}$.
The flow-matching loss is then
\begin{equation}
  \LFM(\theta)
  \;=\;
  \E_{t, x_0, (x_1, \mathbf{C})}
  \bigl\lVert v_\theta(t, \mathbf{C}, x_t) - (x_1 - x_0) \bigr\rVert^{2}.
  \label{eq:fm_loss}
\end{equation}
Crucially, $\LFM$ is a \emph{token-wise} regression: the squared
error is averaged independently over the $N$ audio tokens, and the
target $x_1 - x_0$ for each token depends only on that token's
noise sample. Cross-modal information enters exclusively through
$\mathbf{C}$, and only via the learned velocity field. Whatever
relational structure aligns the audio tokens to the video tokens
must therefore be absorbed by $v_\theta$'s parameters --- there is
no term in \eqref{eq:fm_loss} that explicitly compares the geometry
of the two streams.

\paragraph{The MMAudio backbone.}
We instantiate the velocity field $v_\theta$ via the open-source
\textbf{MMAudio} system of \citet{cheng_2025_mmaudio}, which at the
time of writing is the state of the art among public V2A models trained on
VGGSound and surpasses earlier flow-matching V2A systems such as
Frieren \citep{wang2025frierenefficientvideotoaudiogeneration} (channel-level fusion +
reflow). Three of its design choices are
load-bearing for our methodology and worth describing in full.

\emph{Input encoders.} The conditioning tuple
$\mathbf{C} = (F_v, F_{\mathrm{syn}}, F_t)$ is produced by three
frozen pretrained encoders: a CLIP visual encoder extracts
$F_v \in \R^{n_v \times d_v}$ at $8$ fps ($d_v = 1024$);
Synchformer (a self-supervised audio--visual desynchronization
detector) extracts $F_{\mathrm{syn}}$ at $24$ fps
($d_{\mathrm{syn}} = 768$); and the CLIP text encoder produces
$F_t \in \R^{77 \times d_t}$. Audio is encoded by a pretrained
spectrogram VAE into latents $x_1 \in \R^{N \times d_a}$ at
$31.25$ fps with $d_a = 20$ (S-16kHz model). Because the encoders
are frozen, the only learnable parameters that touch the raw
features are the input projections $\pi_v, \pi_{\mathrm{syn}},
\pi_t, \pi_a$ and the conditioning projector $\pi_g$ that produces
the global vector $c_g$ from the pooled visual / text streams.
This matters for the regularizer of \S\ref{sec:methodology:reps},
which acts precisely at these projection layers.

\emph{Backbone architecture.} The flow predictor is a stack of
$N_1$ {multimodal MM-DiT blocks} adapted from the
text-to-image rectified-flow transformers of
\citet{esser_2024_scaling}, in which separate weights process the
visual / text / audio streams while a single joint
attention operation lets the three streams exchange information,
followed by $N_2$ audio-only transformer blocks. We use the
default S-16kHz configuration ($N_1 = 4$, $N_2 = 8$, hidden
dimension $h = 448$, $157$M trainable parameters). Cross-modal
information is injected at two granularities, and these are the
two carriers our regularizer can target:
\begin{enumerate}
  \item a \emph{global} adaLN modulation by a vector
    $c_g \in \R^{1 \times h}$ pooled from the projected visual
    and text features and broadcast to every audio token, applied
    in every block; and
  \item a \emph{token-level} adaLN modulation by a frame-aligned
    vector $c_f \in \R^{N \times h}$ obtained by upsampling the
    projected Synchformer features to the audio frame rate, applied
    only in the audio stream and intended specifically for
    audio--visual synchronisation.
\end{enumerate}
The global vector $c_g$ is the natural target for a
\emph{semantic} relational signal (one vector per sample, summarising
``what is in the clip''), while the token-level $c_f$ is the
natural target for a \emph{temporal} synchronisation signal. Our
global-conditioning variant in \S\ref{sec:methodology:reps} aligns the
audio batch geometry to the geometry of $c_g$ across samples;
$c_f$ is left untouched because its role is already strongly
constrained by Synchformer.

\emph{Training recipe.} MMAudio is trained jointly on
audio-visual-text data (VGGSound, the only tri-modal source) and
audio-text data (AudioCaps, Clotho, WavCaps), with missing
modalities filled by learnable empty-token embeddings $\varnothing_v,
\varnothing_{\mathrm{syn}}, \varnothing_t$. A classifier-free
guidance schedule masks out visual or text tokens with $10\%$
probability during training. We inherit this recipe unchanged ---
including the joint-training data mix, the empty-token mechanism,
the $300\text{K}$-step training horizon, the AdamW optimiser with
$\beta_2 = 0.95$ (rather than the default $0.999$, which causes
occasional NaNs in their setting), and the linear-then-stepped
learning-rate schedule. Two implications follow for the
regularizer: (i) the GW penalty must be computed only over the
subset of the batch with real video features, since the
empty-token embeddings carry no relational information
(\S\ref{sec:methodology:training}); and (ii) any improvement we
report is attributable to the regularizer alone, not to a change
in the data mix or training schedule.

None of the conditioning mechanisms above align the pairwise
relations among audio tokens to those among video tokens or among
samples in the batch; each token of one stream is nudged toward
the correct token of the other stream by attention and adaLN, but
the geometry of the streams is left to emerge incidentally from
the token-wise regression in \eqref{eq:fm_loss}. The regularizer
below targets exactly this gap.

In our experiments, we focus only on audio-visual regularization through the Gromov-Wasserstein term as the text conditioning signal only as 

\subsection{Representation variants and the relational signal}
\label{sec:methodology:reps}

For a mini-batch of size $B$, write
$F_v^{(b)} \in \R^{n_v \times d_v}$ for the raw CLIP feature
sequence of sample $b$ (with $n_v$ visual tokens of dimension
$d_v = 1024$) and $x_1^{(b)} \in \R^{N \times d_a}$ for the
corresponding audio latent. The regularizer takes as input two
families of squared-Euclidean intra-modal distance matrices, one
per side. We expose five constructions of these matrices, each
corresponding to a different geometric hypothesis about
\emph{which} representations carry the cross-modal relational
signal and at \emph{which} learnable layer it should be injected.
The constructions partition naturally into two families. The first
four operate at the \emph{batch} level: they compress each token
sequence into a single vector per sample and compute pairwise
distances across the $B$ samples of the mini-batch, yielding
$B \times B$ matrices $\DV, \DA$ that compare \emph{clips} to each
other. The fifth operates at the \emph{token} level: it keeps the
token sequences intact and computes pairwise distances within
each clip, yielding per-sample $n \times n$ and $N \times N$
matrices that compare visual and audio tokens of the same clip
to each other.

\paragraph{Encoder-space variant.}
The first construction, which we call the \emph{encoder-space}
variant, computes pairwise distances over pre-projection averages
$\bar v^{(b)} = \tfrac{1}{n_v} \sum_i F_v^{(b),i}$ and
$\bar x^{(b)} = \tfrac{1}{N} \sum_i x_1^{(b),i}$, where the
average is taken over the token dimension. No learnable parameter
sits between the encoders and the GW solver: the gradient of the
regularizer cannot reshape any layer of MMAudio, since the only
parameters in the path are the frozen CLIP and VAE weights. This
variant therefore answers a diagnostic question --- whether the
raw encoder geometry already carries relational structure that
the entropic-GW solver can exploit --- and acts as a control: any
improvement it brings can be attributed entirely to the encoder
geometry as inherited from CLIP, not to representation learning
in the backbone.

\paragraph{Projection-space variant.}
The \emph{projection-space} variant passes the raw token sequences
through MMAudio's input projections $\pi_v, \pi_a$ before
averaging,
$\tilde v^{(b)} = \tfrac{1}{n_v} \sum_i \pi_v(F_v^{(b)})_i$,
analogously for audio. The pairwise-distance matrices are now
computed on representations whose first layer is learnable, so
the GW gradient back-propagates into $\pi_v$ and $\pi_a$ and can
reshape the very inputs the multimodal transformer blocks see.
Because $\pi_a$ is also on the path of every flow-matching
gradient, the GW signal here competes most directly with the
flow-matching loss for control of the projection geometry; this
makes the projection-space variant the most architecturally
``invasive'' of the batch-level alternatives.

\paragraph{Global-conditioning variant.}
The \emph{global-conditioning} variant aligns the audio side to
the very vector that drives MMAudio's adaLN modulation. The audio
representation is computed exactly as in the projection-space
variant, but the video representation is replaced by the
post-projection global conditioning vector
$c_g^{(b)} = \pi_g\bigl(\tfrac{1}{n_v} \sum_i \pi_v(F_v^{(b)})_i\bigr)$,
where $\pi_g$ is the conditioning projector that produces the
adaLN scales and biases broadcast to every token of every block.
Because $c_g$ is the only carrier of \emph{semantic} (as opposed
to temporal-synchronisation) cross-modal information in MMAudio,
this variant directly targets the channel that is already
responsible for ``what is in the clip''. Geometrically, it asks
that the batch geometry of the conditioning summary match the
batch geometry of the audio targets, allowing the regularizer to
shape both $\pi_a$ and the conditioning projector $\pi_g$.

\paragraph{Fused variant.}
The \emph{fused} variant uses the same projected representations
as the projection-space variant but couples the relational GW
objective to a Wasserstein term on a cross-modal cosine cost,
yielding a fused Gromov--Wasserstein loss in the sense of
\citet{vayer_2019_optimal, vayer_2019_fused}. The two terms share
a single Sinkhorn-solved transport plan, following the GOT recipe
of \citet{chen_2020_graph}. The cross-modal cost is constructed
exactly as described in \S\ref{sec:methodology:fgw} below: each
sample's projected video and audio representations are
$L_2$-normalised onto the unit sphere, and the cross-modal cost
matrix has entries $1 - \hat v_i^{\!\top} \hat a_j$. The diagonal
of this matrix records the cosine distance between the projected
video and audio representations of the \emph{same} sample, which
the Wasserstein term penalises directly while the GW term
continues to align the relational geometry of the two streams.

\paragraph{Token-graph variant.}
The fifth construction, which we call the \emph{token-graph}
variant, abandons the cross-batch view and instead solves a fused
GW problem \emph{per sample, on token-level graphs}. For each
sample $b$ in the mini-batch we form an intra-video distance
matrix
$\DV^{(b)} \in \R^{n_v \times n_v}$
from the projected visual tokens, an intra-audio distance matrix
$\DA^{(b)} \in \R^{N \times N}$
from the projected audio tokens, and a per-sample cross-modal
cosine cost
$C_{\mathrm{xy}}^{(b)} \in \R^{n_v \times N}$
between every visual token and every audio token of the same
clip. A fused GW problem with a single shared $n_v \times N$
transport plan is solved independently per sample, and the loss
is averaged over the batch. This is the formulation closest in
spirit to the original GOT
\citep{chen_2020_graph}, lifted from image-region / word-token
sets to within-clip video / audio token sets. Because each sample
contributes its own transport plan and its own pair of distance
matrices, the token-graph variant cannot collapse the batch onto
a single direction in the way the four batch-level variants can,
and consequently does not require the spectral anti-collapse
penalty introduced in \S\ref{sec:methodology:anticollapse}.

Table~\ref{tab:variants} summarises the five constructions along
the dimensions that distinguish them: which representation feeds
the distance matrices, which learnable parameters lie on the GW
gradient path, the shape of the resulting cost matrices, and
whether a cross-modal Wasserstein term is added.

\begin{table}[t]
  \centering
  \small
  \begin{tabular}{@{}lllll@{}}
    \toprule
    Variant & Representation & Learnable parameters on GW path & Distance-matrix shape & Cross-modal cost? \\
    \midrule
    Encoder-space        & $\bar v$, $\bar x$                                 & none                             & $B \times B$                              & no  \\
    Projection-space     & $\pi_v(F_v)$, $\pi_a(x_1)$ averaged                & $\pi_v$, $\pi_a$                 & $B \times B$                              & no  \\
    Global-conditioning  & $c_g = \pi_g(\overline{\pi_v F_v})$, $\pi_a(x_1)$  & $\pi_v$, $\pi_g$, $\pi_a$        & $B \times B$                              & no  \\
    Fused                & $\pi_v(F_v)$, $\pi_a(x_1)$ averaged                & $\pi_v$, $\pi_a$                 & $B \times B$                              & yes (batch-level) \\
    Token-graph          & $\pi_v(F_v)$, $\pi_a(x_1)$ token-wise              & $\pi_v$, $\pi_a$                 & $n_v \times n_v$, $N \times N$ per sample & yes (per-sample) \\
    \bottomrule
  \end{tabular}
  \caption{The five representation variants used by the GW
    regularizer, ordered from least to most architecturally
    invasive on the four batch-level rows. For all variants
    except encoder-space, the gradient of the regularizer reaches
    at least the input projections of MMAudio; the
    global-conditioning variant additionally reshapes the
    conditioning projector $\pi_g$ that drives every adaLN block.
    The token-graph variant is the only one whose distance
    matrices are intra-clip rather than inter-clip; it is also
    the only one that does not require the spectral
    anti-collapse penalty of
    \S\ref{sec:methodology:anticollapse}.}
  \label{tab:variants}
\end{table}

For all variants, raw squared distances $D$ are normalised by
their mean,
\begin{equation}
  \widehat D \;=\; D \,/\, \overline{D},
  \qquad
  \overline{D} \;=\; \mathrm{mean}(D),
  \label{eq:dist_normalisation}
\end{equation}
with the divisor detached from the autograd graph. This makes the
regularizer scale-invariant in the sense that multiplying any
representation by a constant leaves $\widehat D$ unchanged, which in
turn makes the trade-off coefficient $\lambda_{\mathrm{GW}}$
introduced in \S\ref{sec:methodology:training} comparable across
variants and across training steps even as feature magnitudes
drift.

\subsection{Entropic Gromov--Wasserstein}
\label{sec:methodology:egw}

Let $\widehat\DV, \widehat\DA \in \R^{B \times B}$ denote the
normalised intra-video and intra-audio distance matrices for one of
the cross-batch variants above, and equip both sides with the
uniform marginal
$p = q = \tfrac{1}{B}\mathbf{1}_B$. The quadratic
Gromov--Wasserstein discrepancy with squared cost is
\begin{equation}
  \mathrm{GW}^{2}(\widehat\DV, \widehat\DA, p, q)
  \;=\;
  \min_{T \in \Pi(p,q)}
  \sum_{i,j,k,l}
  \bigl(\widehat\DV_{ik} - \widehat\DA_{jl}\bigr)^{2}
  \, T_{ij} T_{kl},
  \label{eq:gw}
\end{equation}
where $\Pi(p,q) = \{T \in \R_{\geq 0}^{B \times B} :
T \mathbf{1}_B = p,\, T^{\!\top}\!\mathbf{1}_B = q\}$. Equation
\eqref{eq:gw} is a quadratic assignment problem and NP-hard in
general; following \citet{peyre_2016_gromov} and
\citet{solomon_2016_entropic} we work with its entropic relaxation
\begin{equation}
  \mathrm{EGW}_\varepsilon
  \;=\;
  \min_{T \in \Pi(p,q)}
  \;\bigl\langle L(\widehat\DV, \widehat\DA) \otimes T,\, T\bigr\rangle
  \;-\; \varepsilon\, H(T),
  \label{eq:egw}
\end{equation}
with $L(\widehat\DV, \widehat\DA)_{ijkl} = (\widehat\DV_{ik} -
\widehat\DA_{jl})^{2}$ and $H(T) = -\sum_{ij} T_{ij}(\log T_{ij} -
1)$ the discrete entropy. The entropic relaxation is convex in $T$
\emph{conditional on the linearisation around the current iterate},
and admits the Sinkhorn-iteration solver below; its statistical
behaviour at finite $B$ is governed by the parametric
$B^{-1/2}$ rate of \citet{zhang_2024_gromov} and the convergence
guarantees of \citet{rioux_2023_entropic}, neither of which holds
for the un-regularised problem.

\paragraph{Mirror-descent Sinkhorn solver.}
The entropic-GW solver follows the mirror-descent / Frank--Wolfe
recipe of \citet{peyre_2016_gromov} (introduced for GW barycenter
computation) and \citet{solomon_2016_entropic} (provably
convergent on entropic-GW correspondence problems). The idea is
to alternate two cheap steps: a \emph{linearisation} of the
quartic GW interaction around the current plan, and a single
\emph{entropic optimal-transport projection} onto the marginal
polytope solved by Sinkhorn iteration.

For squared cost the GW objective in \eqref{eq:gw} admits the
linearised gradient
\begin{equation}
  G(T) \;=\; -2\, \widehat\DV \, T \, \widehat\DA
  \;\;+\;\; \text{constants independent of } T,
  \label{eq:gw_gradient}
\end{equation}
since
$\partial / \partial T \bigl\langle L(\widehat\DV, \widehat\DA)
\otimes T, T\bigr\rangle = 2 (\widehat\DV^{2} p \mathbf{1}^{\!\top}
+ \mathbf{1} q^{\!\top} \widehat\DA^{2}) - 2 \widehat\DV T
\widehat\DA$, and the marginal-dependent terms are constant under
fixed $p, q$. At each outer iteration we therefore solve an
entropic optimal transport problem with cost matrix $G(T)$, that is,
\begin{equation}
  T^{+}
  \;=\;
  \arg\min_{T' \in \Pi(p,q)}
  \langle G(T), T'\rangle - \varepsilon H(T'),
  \label{eq:inner_eot}
\end{equation}
whose solution has the Gibbs form $T^{+}_{ij} = u_i K_{ij} v_j$
with $K = \exp(-G(T)/\varepsilon)$, and $u, v$ are recovered by the
Sinkhorn fixed-point recursion
\begin{equation}
  u \leftarrow p \,/\, (K v),
  \qquad
  v \leftarrow q \,/\, (K^{\!\top} u),
  \label{eq:sinkhorn}
\end{equation}
iterated until the marginals match to a tolerance fixed by the
inner iteration count. The recursion is the standard Bregman /
matrix-scaling iteration that defines the entropic-OT solution
and inherits its linear convergence in the marginal-error
metric; we treat its hyperparameters
(\texttt{inner\_sinkhorn\_iter}, $\varepsilon$) as fixed budgets
rather than tuning to a stopping criterion. Convergence of the
overall outer-inner scheme on entropic GW with quadratic cost is
guaranteed by \citet{rioux_2023_entropic}, who prove
that the entropic-GW objective is convex in the linearised
regime and admits provable rates --- the theoretical underpinning
that justifies running only a small number of outer steps in
practice.

We run \texttt{num\_sinkhorn\_iter}$=5$ outer mirror-descent
steps, each containing \texttt{inner\_sinkhorn\_iter}$=20$
Sinkhorn passes, with $\varepsilon = 0.1$, and we stabilise the
kernel with the standard log-sum-exp shift
$K = \exp(-G/\varepsilon - \max(-G/\varepsilon))$ that
\citet{feydy_2019_interpolating} popularised in their
GPU-friendly Sinkhorn implementation. The full computation is
wrapped in an \texttt{autocast(enabled=False)} block so that the
entire solver runs in \texttt{float32}, which we found necessary
to avoid overflow in $K$ when $\varepsilon$ is small or distances
are poorly scaled --- a known sensitivity of entropic-OT solvers
addressed in the same line of work
\citep{feydy_2019_interpolating, rioux_2023_entropic}.

\paragraph{Loss value.}
At the converged plan $T^{\star}$ we report the full quadratic
objective rather than the linearised gradient. With uniform
marginals the squared-cost expansion reads
\begin{equation}
  \LGW
  \;=\;
  \langle \widehat\DV^{\odot 2}, p p^{\!\top} \rangle
  \;+\;
  \langle \widehat\DA^{\odot 2}, q q^{\!\top} \rangle
  \;-\;
  2\,\langle \widehat\DV\, T^{\star}\, \widehat\DA,\, T^{\star} \rangle,
  \label{eq:gw_value}
\end{equation}
where $D^{\odot 2}$ denotes element-wise square. The first two
terms are constants in $T^{\star}$ but vary across mini-batches and
across training steps; reporting them keeps loss magnitudes
comparable across the four variants.

\subsection{Fused Gromov--Wasserstein}
\label{sec:methodology:fgw}

For the fused and token-graph variants we add a
linear (Wasserstein) cross-modal term and share a single transport
plan between the two contributions, following GOT
\citep{chen_2020_graph} and the FGW formulation of
\citet{vayer_2019_optimal, vayer_2019_fused}.

\paragraph{Cross-modal cosine cost.}
The pure-GW objective \eqref{eq:gw_value} compares
\emph{intra-modal pairwise distances} only; it is invariant to any
isometric reparametrisation of either modality, including
arbitrary global rotations or reflections. The fused term breaks
this invariance by re-introducing a direct video-to-audio
\emph{pointwise} comparison through the cross-modal cost matrix
$C_{\mathrm{xy}}$. Concretely, write
$\tilde v, \tilde a \in \R^{\bullet \times d}$ for the projected
representations on the video and audio sides --- where $\bullet$
is the batch index $b \in \{1,\dots,B\}$ for the cross-batch
fused variant, and the within-sample token index for the
token-graph variant. We first $L_2$-normalise each row so that
every representation lies on the unit sphere of $\R^d$,
\begin{equation}
  \hat v_i \;=\; \tilde v_i \,/\, \lVert \tilde v_i \rVert_2,
  \qquad
  \hat a_j \;=\; \tilde a_j \,/\, \lVert \tilde a_j \rVert_2,
  \label{eq:l2norm}
\end{equation}
and then compute the matrix of pairwise cosine \emph{distances}
\begin{equation}
  \bigl[C_{\mathrm{xy}}\bigr]_{ij}
  \;=\;
  1 \;-\; \hat v_i^{\!\top} \hat a_j
  \;\in\; [0, 2].
  \label{eq:ccross}
\end{equation}
For the cross-batch fused variant this yields
$C_{\mathrm{xy}} \in \R^{B \times B}$ whose diagonal entries are
the cosine distances between the projected video and audio
representations of the \emph{same} sample $b$ --- the natural
"pointwise" cost that an OT plan with diagonal mass would pay zero
under perfect cross-modal feature alignment. For
the token-graph variant the same construction is performed
\emph{within each sample} on the projected token sequences,
yielding a per-sample $C_{\mathrm{xy}}^{(b)} \in \R^{n \times m}$
that compares each video token to each audio token of the same
clip. Three choices deserve note. (i) The cosine form is preferred
to a Euclidean distance because $\tilde v$ and $\tilde a$ live in
different feature spaces with potentially very different scales:
normalising to the unit sphere removes the scale ambiguity and
keeps $C_{\mathrm{xy}} \in [0, 2]$ for any $d$ and any feature
magnitude, so $C_{\mathrm{xy}}$ and the mean-normalised intra-modal
costs $\widehat\DV, \widehat\DA$ live on comparable scales without
further rescaling. (ii) The construction reuses the projection
layers $\pi_v, \pi_a$ that already produce the intra-modal
representations, so the FGW objective adds no new learnable
parameters: the same projector outputs feed the GW and Wasserstein
terms, ensuring they pull on the same representation rather than on
two disjoint heads. (iii) The video side of $C_{\mathrm{xy}}$
inherits the asymmetric \texttt{detach\_video} treatment of
\S\ref{sec:methodology:training}, so the Wasserstein term too acts
as a one-way pressure on the audio representation.

\paragraph{Linearised gradient and loss.}
With $C_{\mathrm{xy}}$ in hand and mixing parameter
$\alpha \in [0,1]$, the linearised gradient at iterate $T$ becomes
\begin{equation}
  G_{\mathrm{FGW}}(T)
  \;=\;
  \alpha \cdot \bigl(-2\, \widehat\DV \, T \, \widehat\DA\bigr)
  \;+\;
  (1 - \alpha) \cdot C_{\mathrm{xy}},
  \label{eq:fgw_gradient}
\end{equation}
plugged into the same inner entropic-OT solve
\eqref{eq:inner_eot}--\eqref{eq:sinkhorn}. The reported loss is
the same convex combination evaluated at $T^{\star}$,
\begin{equation}
  \LFGW
  \;=\;
  \alpha \cdot \LGW(T^{\star})
  \;+\;
  (1 - \alpha) \cdot \langle C_{\mathrm{xy}},\, T^{\star} \rangle.
  \label{eq:fgw_value}
\end{equation}
\textit{Sharing} the plan across the two terms is non-trivial: it
forces the Wasserstein and Gromov--Wasserstein contributions to
agree on a single coupling, rather than letting each pull the
representations toward its own optimum. \citet{chen_2020_graph}
report empirically that the shared-plan variant outperforms two
independent plans on every cross-domain task they examine; we
adopt the same design.

\subsection{Anti-collapse penalty}
\label{sec:methodology:anticollapse}

A practical risk of pure GW alignment between cross-batch
representations is that the projector(s) collapse the batch onto
a single direction, sending $\widehat\DV$ and $\widehat\DA$ both
toward zero and trivially minimising \eqref{eq:gw_value}. This is
a structural pathology of any similarity-driven objective with a
learnable feature extractor and no negative-pair pressure --- the
same failure mode that motivated the design of non-contrastive
self-supervised methods such as Barlow Twins
\citep{zbontar_2021_barlow} (cross-correlation matrix toward
identity), VICReg \citep{bardes_2022_vicreg} (per-dimension
variance and off-diagonal covariance penalties), and W-MSE
\citep{ermolov_2021_whitening} (hard whitening of batch
embeddings). At a higher level it also resembles the
entropic-bias collapse of plain entropic-OT losses that motivated
the Sinkhorn-divergence construction of
\citet{feydy_2019_interpolating}. We adopt the **spectral
log-determinant** form of the regularizer, introduced by
\citet{ozsoy_2022_corinfomax} as the second-order
mutual-information surrogate underlying the CorInfoMax objective.
Concretely, we add a spectral penalty on the centred Gram matrix
\begin{equation}
  K(v) \;=\; \frac{1}{B-1}\,\bigl(v - \bar v\bigr)\bigl(v - \bar v\bigr)^{\!\top}
  \;\in\; \R^{B \times B},
  \label{eq:gram}
\end{equation}
where $v \in \R^{B \times d}$ stacks the batch representations,
and define
\begin{equation}
  \mathcal{L}_{\mathrm{AC}}(v)
  \;=\;
  -\log\det\bigl(K(v) + \eta I_B\bigr).
  \label{eq:anticollapse}
\end{equation}
The non-zero spectrum of the Gram $K(v)$ coincides with that of the
$d \times d$ feature covariance up to the $(d - B)\log\eta$
constant when $d > B$, so the Gram form is computationally
preferable when $d \gg B$ but penalises the same rank deficiency
as the covariance form of \citet{ozsoy_2022_corinfomax}.
A rank-1 collapse drives $B-1$ eigenvalues of $K$ to zero, sending
\eqref{eq:anticollapse} toward $(B-1)\,(-\log\eta)$ — a large
positive number that repels the optimiser. Compared to VICReg's
per-dimension variance hinge or Barlow Twins' off-diagonal squared
penalty, the log-det form is a \emph{soft barrier on the joint
spectrum} rather than a per-axis or per-pair constraint: it forbids
rank deficiency without dictating the anisotropy of the surviving
spectrum, which leaves the GW objective free to choose the
direction of variation. This is the same trade-off that
distinguishes soft whitening (CorInfoMax-style) from hard
whitening (W-MSE-style \citep{ermolov_2021_whitening}); we prefer
the soft form because the GW geometry is precisely what we want to
shape, not constrain. Critically, the penalty is applied to the
\emph{non-detached} representations, so its gradient always reaches
the projector even when the GW term is computed with the video
side detached (\S\ref{sec:methodology:training}).

\subsection{Training integration and the \texorpdfstring{$\lambda_{\mathrm{GW}}$}{lambda\_GW} schedule}
\label{sec:methodology:training}

The total loss optimised at training step $s$ is
\begin{equation}
  \mathcal{L}_{\mathrm{total}}(\theta; s)
  \;=\;
  \LFM(\theta)
  \;+\;
  \lambda_{\mathrm{GW}}(s) \cdot \mathcal{L}_{\mathrm{rel}}(\theta)
  \;+\;
  \lambda_{\mathrm{AC}} \cdot \mathcal{L}_{\mathrm{AC}}(\theta),
  \label{eq:total_loss}
\end{equation}
where $\mathcal{L}_{\mathrm{rel}} \in \{\LGW, \LFGW\}$ is selected
by the variant. Three implementation choices deserve emphasis.

\paragraph{Subset selection.}
The regularizer is computed only over the subset of the mini-batch
for which a real video is present (\texttt{video\_exist} mask),
since for audio-only samples MMAudio's own training procedure
replaces the visual stream with learnable empty tokens; aligning
audio geometry to the geometry of constants is meaningless. If
fewer than \texttt{min\_batch}$=4$ real videos remain in the
batch, the GW term is set to zero for that step.

\paragraph{Asymmetric gradient flow.}
By default we set \texttt{detach\_video}$=\mathtt{True}$, which
stops the gradient on the video side of $\widehat\DV$ before the GW
solver. The video encoder (CLIP / Synchformer) is frozen in any
case, but the projection layers $\pi_v$ and $\pi_g$ are not, and
without the detach the GW loss can satisfy itself by reshaping the
visual conditioning rather than the audio representation. Detaching
ensures that the GW term acts as a one-way pressure: the audio side
is asked to mirror the video geometry, not the other way round.
The anti-collapse penalty bypasses this detach so it can still
discipline the projector outputs.

\paragraph{Schedule.}
The coefficient $\lambda_{\mathrm{GW}}(s)$ is governed by a small
schedule family:
\begin{equation}
  \lambda_{\mathrm{GW}}(s)
  \;=\;
  \begin{cases}
    \lambda_{0} \cdot s / s_{\mathrm{w}}
      & \text{if } s < s_{\mathrm{w}} \;\text{and schedule}=\textsc{linear-rampup}, \\[2pt]
    0
      & \text{if } s < s_{\mathrm{w}} \;\text{and schedule}=\textsc{constant}, \\[2pt]
    \lambda_{0}
      & \text{if } s \geq s_{\mathrm{w}} \;\text{and schedule}\in\{\textsc{constant}, \textsc{linear-rampup}\}, \\[2pt]
    \tfrac{\lambda_{0}}{2}\bigl(1 + \cos\bigl(\pi\,\tfrac{s - s_{\mathrm{w}}}{S - s_{\mathrm{w}}}\bigr)\bigr)
      & \text{if } s \geq s_{\mathrm{w}} \;\text{and schedule}=\textsc{cosine-anneal},
  \end{cases}
  \label{eq:lambda_schedule}
\end{equation}
with base weight $\lambda_{0}$, warm-up length $s_{\mathrm{w}}$,
and total training horizon $S$. The motivation is empirical and
theoretical at once. Empirically: at initialisation the audio
representation has not yet been shaped by the flow-matching loss,
so its intra-batch geometry is essentially noise, and forcing it
to match the video geometry from step zero pulls the network into
degenerate solutions that satisfy the GW term at the cost of the
generative objective. Theoretically: existing fused-OT work (e.g.\
\citet{chen_2020_graph}) fixes the mixing $\alpha$ but not its
relative weight against the task loss. The schedule
\eqref{eq:lambda_schedule} resolves the latter dimension explicitly
and is, to our knowledge, novel for OT-based cross-modal
regularization.

\subsection{Why this isolates the relational hypothesis}
\label{sec:methodology:isolation}

Three properties make the construction above a clean test of the
thesis question. First, $\LFM$ in \eqref{eq:fm_loss} and the
backbone of \citet{cheng_2025_mmaudio} together constitute a
\emph{pointwise} alignment baseline: the generative target is
defined token by token via the constant velocity $x_1 - x_0$, and
all cross-modal information enters via mechanisms (joint attention,
adaLN with $c_g$, frame-aligned $c_f$) that operate on individual
tokens rather than on inter-token relations. The GW term
\eqref{eq:gw_value} is the smallest principled addition that
introduces an explicit \emph{relational} signal, because
Gromov--Wasserstein is, by construction, the OT object that
compares pairwise distance matrices rather than features.
Improvements traceable to $\LGW$ can therefore be attributed to
the relational signal specifically, rather than to extra capacity
or extra supervision in some other form.

Second, the five representation variants of
\S\ref{sec:methodology:reps} disentangle \emph{where} in the
network the relational signal is injected. The encoder-space
variant operates before any learnable layer and asks whether the
encoder geometry already supports relational alignment without
parameter updates; the projection-space variant injects the
signal into the input projections and lets it propagate through
every transformer block; the global-conditioning variant aligns
the very vector that drives adaLN modulation; the fused variant
adds a Wasserstein term to test whether relational and pointwise
signals are complementary at the same representation level; and
the token-graph variant moves the GW objective from cross-batch
geometry to within-sample token geometry, matching most closely
the setting in which GW is typically deployed
\citep{chen_2020_graph, vayer_2019_optimal}. The five variants
are the five most natural injection points in MMAudio's
information flow, and the ablation is exhaustive at this
granularity.

Third, the schedule \eqref{eq:lambda_schedule} resolves a confound
that prior fused-OT work leaves implicit: improvements (or
regressions) cannot be attributed to a fortuitous choice of
$\lambda$ if the same schedule is used across all variants, and
ablating the schedule itself separates relational signal from
training dynamics.

The contribution of this chapter, relative to the upstream
MMAudio repository, is therefore localised in three additions: a
self-contained GW solver
(\texttt{src/mmaudio/model/gw\_regularization.py}), the integration
hook in the training loop that selects the active subset, applies
the schedule, and combines the two losses, and the configuration
surface (\texttt{config/base\_config.yaml::gw\_regularization}) that
exposes the variant, the schedule, the entropic regularization
$\varepsilon$, the FGW mixing $\alpha$, the detach flag, and the
anti-collapse parameters as the axes of the empirical study in
\S\ref{sec:results}. The flow-matching backbone, the encoders, the
synchronization module, the data pipeline, and every other
architectural choice are inherited unchanged from
\citet{cheng_2025_mmaudio}; this isolates the regularizer as the
single independent variable.
